% =============================================================================
% Defensa del trabajo final. Acompaña a main.tex y no lo reemplaza: cada cifra
% de estas diapositivas debe existir primero en la tesis y en evidence/.
%
% Compilar:
%   pdflatex --enable-installer -interaction=nonstopmode presentacion.tex
%   pdflatex --enable-installer -interaction=nonstopmode presentacion.tex
% =============================================================================
\documentclass[10pt,aspectratio=169]{beamer}

\usepackage[utf8]{inputenc}
\usepackage[T1]{fontenc}
\usepackage[spanish,es-nodecimaldot]{babel}
\usepackage{lmodern}
\usepackage{booktabs}
\usepackage{amsmath,amssymb}

% Version con notas del orador (pantalla dual):
%   \setbeameroption{show notes on second screen=right}
% Para imprimir solo las notas:
%   \setbeameroption{show only notes}
\usepackage{pgfpages}
\usetheme{default}
\usecolortheme{seahorse}
\setbeamertemplate{navigation symbols}{}
\setbeamertemplate{footline}[frame number]
\setbeamertemplate{itemize items}[circle]

\title[ManyHands]{Orquestación de agentes de lenguaje para el desarrollo de software}
\subtitle{Una política de descomposición adaptativa con verificación basada en evidencia}
\author{Francisco Clarke \\ \small LU 120547}
\institute{\small Departamento de Ciencias e Ingeniería de la Computación \\ Universidad Nacional del Sur \\[0.5em] \scriptsize Directora: Dra. Ana Gabriela Maguitman \\ \scriptsize Codirector: Dr. Federico Martín Schmidt}
\date{2026}

\begin{document}

\begin{frame}[plain]
  \titlepage
\end{frame}

% -----------------------------------------------------------------------------
\section{Problema}

\begin{frame}{El problema}
  Un agente de lenguaje puede resolver una tarea de programación acotada. Un
  objetivo real de software no es una tarea acotada.

  \bigskip
  \pause
  La salida natural es \textbf{descomponer}. Y ahí aparece la tensión:

  \bigskip
  \begin{columns}[T]
    \begin{column}{0.47\textwidth}
      \textbf{Dividir de menos}
      \begin{itemize}
        \item El agente satura su contexto.
        \item No hay puntos de verificación intermedios.
        \item Un fallo invalida todo el trabajo.
      \end{itemize}
    \end{column}
    \begin{column}{0.47\textwidth}
      \textbf{Dividir de más}
      \begin{itemize}
        \item Aparecen contratos entre unidades.
        \item Crece el costo de integración.
        \item Los conflictos serializan el trabajo.
      \end{itemize}
    \end{column}
  \end{columns}

  \bigskip
  \pause
  \begin{center}
    \textbf{La paradoja de la granularidad:} el óptimo no es un extremo, y no es
    el mismo para todas las tareas.
  \end{center}
\note{%
  Abrir por la tension, no por la herramienta. El tribunal tiene que entender
  que ninguno de los dos extremos es una solucion: dividir de menos satura al
  agente, dividir de mas traslada el costo a la integracion. Si preguntan por
  que no basta con dividir bien a mano, la respuesta es que el optimo depende
  de la tarea, y por eso hace falta una politica.%
}
\end{frame}

\begin{frame}{Preguntas de investigación}
  \begin{description}
    \item[RQ1] ¿Cómo varía la tasa de entrega verificada entre hoja única
      forzada, división fina fija y una política adaptativa?
    \item[RQ2] ¿Qué trade-off existe entre éxito, duración, costo y overhead de
      coordinación?
    \item[RQ3] ¿Qué modos de falla aparecen en cada configuración?
  \end{description}

  \bigskip
  \begin{center}
    \small Estudio \textbf{exploratorio} de viabilidad y modos de falla.\\
    No busca —ni podría alcanzar— significancia estadística.
  \end{center}
\note{%
  Decir explicitamente que es exploratorio ANTES de mostrar cualquier numero.
  Si alguien pide significancia, la respuesta ya esta dada: con dos
  observaciones por celda no hay base para un p-valor, y afirmar lo contrario
  seria deshonesto.%
}
\end{frame}

% -----------------------------------------------------------------------------
\section{Aporte}

\begin{frame}{La política de descomposición adaptativa}
  Cada unidad recibe cuatro señales normalizadas en $[0,10]$:

  \[
    C_{\mathit{task}} = w_1 S_r + w_2 I_i + w_3 V_s + w_4 T_m
  \]

  \begin{center}
  \small
  \begin{tabular}{llc}
    \toprule
    Símbolo & Dimensión & Peso \\
    \midrule
    $S_r$ & Radio de alcance (archivos y módulos afectados) & 0{,}30 \\
    $I_i$ & Impacto sobre interfaces exportadas & 0{,}25 \\
    $V_s$ & Superficie de validación & 0{,}25 \\
    $T_m$ & Masa de contexto en tokens & 0{,}20 \\
    \bottomrule
  \end{tabular}
  \end{center}

  \bigskip
  Una unidad se mantiene como hoja si $C_{\mathit{task}} \le 3{,}5$.

  \bigskip
  \pause
  \textbf{Validación híbrida:} el modelo \emph{propone} las señales; un
  validador determinista las \emph{acota} contra la superficie real del
  repositorio. El origen de cada señal queda registrado como \texttt{llm},
  \texttt{clamped} o \texttt{derived}.
\note{%
  Los pesos NO estan calibrados empiricamente: son una asignacion razonada, y
  eso esta declarado como limitacion en la tesis. Si preguntan por la
  calibracion, reconocerlo de frente y remitir a trabajo futuro. Lo defendible
  aca es la validacion hibrida: el modelo propone, el validador determinista
  acota, y el origen de cada senal queda registrado.%
}
\end{frame}

\begin{frame}{Dónde está la frontera}
  \begin{center}
    \large
    La política decide \textbf{si} dividir.\\[0.6em]
    Solo el Arquitecto puede decir \textbf{cómo}.
  \end{center}

  \bigskip
  \pause
  Esto no es una decisión de escritorio: es un \textbf{resultado negativo}
  obtenido de un run real.

  \begin{itemize}
    \item Una versión previa fabricaba sub-unidades particionando rutas de
      forma mecánica cuando el Arquitecto no proponía ninguna.
    \item En un run canónico los tres agentes terminaron con código de salida
      0, y \textbf{los tres candidatos fueron rechazados por violar su propio
      contrato de alcance}.
    \item Causa: implementar una porción de la API también necesita el tipo del
      dominio del que depende. Una partición de rutas no es una unidad de
      trabajo coherente.
  \end{itemize}

  \pause
  \bigskip
  Corrección: se retiró la síntesis mecánica. Sin sub-unidades propuestas se
  conserva la hoja cohesiva y se registra \texttt{resplit\_declined}.
\note{%
  Esta es la diapositiva mas importante de la defensa. Es un resultado
  negativo, y conviene presentarlo como tal: el sistema hacia algo que parecia
  razonable, un run real lo refuto, y la correccion fue quitar funcionalidad,
  no agregarla. Conecta con Parnas 1972: los modulos se separan por decisiones
  de diseno, no por archivos. Si preguntan como se detecto: tres agentes con
  exit 0 y tres candidatos rechazados por alcance.%
}
\end{frame}

% -----------------------------------------------------------------------------
\section{Arquitectura}

\begin{frame}{Recorrido de un run}
  \begin{enumerate}
    \item \textbf{Grounding} — se indexa el repositorio objetivo; el plan cita
      rutas que existen.
    \item \textbf{Planificación} — el planificador emite unidades semánticas y
      señales de complejidad.
    \item \textbf{Política adaptativa} — $C_{\mathit{task}}$ reforma el árbol;
      la decisión se persiste como evento.
    \item \textbf{Compilación} — grafo con cuatro relaciones tipadas:
      pertenencia, artefactos, costuras y restricciones de conflicto.
    \item \textbf{Ejecución} — cada hoja corre en un worktree de Git aislado,
      con \emph{lease} y \emph{fencing}.
    \item \textbf{Validación} — matriz de evidencias sobre el commit exacto.
    \item \textbf{Integración} — bottom-up, con reparación semántica.
    \item \textbf{Entrega} — manifiesto y receipt sobre el árbol validado.
  \end{enumerate}
\end{frame}

\begin{frame}{Dos invariantes que sostienen todo}
  \textbf{1. \texttt{git diff} es la única fuente de verdad.}\\
  El agente no commitea. El orquestador inspecciona el árbol y commitea el diff
  exacto. Lo que el agente \emph{dice} que hizo no entra en la evidencia.

  \bigskip
  \textbf{2. El alcance es un contrato, no una sugerencia.}\\
  Cada nodo declara las rutas que puede modificar. Un archivo fuera de contrato
  invalida el candidato.

  \bigskip
  \pause
  \small
  Este segundo invariante tuvo un costo medido: bajo alcance estricto, un
  archivo de prueba \emph{nuevo} que el planificador no anticipó hacía fallar
  trabajo correcto. La corrección fue \textbf{creación acotada}: un nodo puede
  crear archivos bajo los directorios que ya posee —roots derivados por el
  compilador, nunca pedidos por el modelo—, pero editar un archivo preexistente
  no declarado sigue fuera de alcance.
\note{%
  Insistir en que el alcance estricto no se relajo para hacer pasar los runs.
  Se acoto la autorizacion a la creacion de archivos nuevos bajo directorios
  que el nodo ya posee, y los roots los deriva el compilador, no el modelo.
  Editar un archivo preexistente no declarado sigue fuera de alcance.%
}
\end{frame}

% -----------------------------------------------------------------------------
\section{Evaluación}

\begin{frame}{Cómo se evaluó}
  \textbf{Caso canónico:} repositorio objetivo externo, commit base congelado,
  objetivo que cruza dominio, API, superficie web y pruebas. Codex CLI
  (\texttt{gpt-5.5}) en planificación y ejecución.

  \bigskip
  \textbf{Estudio comparativo:} 2 tareas $\times$ 3 condiciones $\times$ 2
  repeticiones. Las dos tareas caen \textbf{en lados opuestos del umbral}: si
  ambas fueran complejas, la condición de hoja única fallaría por construcción y
  la política adaptativa parecería superior por una razón trivial.

  \bigskip
  \pause
  Todo se fijó \textbf{antes} de observar resultados: tareas, orden, métricas,
  hipótesis, su falsador y una regla de inconclusión.

  \bigskip
  \small Las tablas y figuras se derivan automáticamente del journal de
  eventos. Ninguna cifra se transcribe a mano.
\note{%
  El punto no obvio de la seleccion de tareas: las dos tareas estan en lados
  opuestos del umbral a proposito. Una politica adaptativa util debe parecerse
  a la hoja unica en la tarea acotada y no ser peor que la division fina en la
  tarea multi-capa. Una politica que siempre divide no resolvio la paradoja:
  eligio un extremo. Con dos tareas complejas eso seria invisible.%
}
\end{frame}

\begin{frame}{El resultado: la hipótesis quedó falsada}
  \small
  \begin{center}
  \begin{tabular}{llccc}
    \toprule
    \textbf{Tarea} & \textbf{Cond.} & \textbf{Entregas} & \textbf{Reloj (s)} & \textbf{Tokens} \\
    \midrule
    T1 multi-capa & A \emph{no dividir} & \textbf{2/2} & 259, 357 & 22k, 25k \\
    T1 multi-capa & B división fina & 1/2 & 942, 1209 & 112k, 101k \\
    T1 multi-capa & C adaptativa & 1/2 & 810, 1109 & 104k, 96k \\
    \midrule
    T2 acotada & A & 2/2 & 282, 291 & 30k, 28k \\
    T2 acotada & B & 2/2 & 362, 393 & 52k, 28k \\
    T2 acotada & C & 2/2 & 344, 351 & 32k, 26k \\
    \bottomrule
  \end{tabular}
  \end{center}

  \bigskip
  \pause
  \textbf{En T2 la hipótesis se sostiene, y por el mecanismo previsto:} las tres
  condiciones convergieron a \emph{una sola unidad} —incluida la que fuerza
  dividir—. Sobre una tarea cohesiva, dividir no es una opción disponible.

  \bigskip
  \pause
  \textbf{En T1 se falsa, en la dirección más desfavorable:} no descomponer
  entregó 2/2, con cerca de un tercio del tiempo y un cuarto de los tokens, y la \emph{misma}
  superficie funcional.
\note{%
  No suavizar esto. El diseno se pre-registro justamente para que un resultado
  desfavorable no pudiera reinterpretarse despues. Si preguntan si entonces el
  aporte sirve: la respuesta honesta es que sobre este objetivo y este
  repositorio no hay evidencia de ventaja, y que el experimento no probo su
  hipotesis donde deberia sostenerse.%
}
\end{frame}

\begin{frame}{Por qué tampoco autoriza la conclusión inversa}
  \begin{itemize}
    \item \textbf{El repositorio es chico} —5 pruebas de base, aprox.\ 25k tokens por
      run—. La descomposición se hipotetiza útil cuando un agente \emph{satura su
      contexto}. Eso no pasó ni cerca: el experimento no puso a prueba su
      hipótesis en el régimen que le es favorable.
    \item \textbf{Dos celdas discrepan entre repeticiones} → la varianza del
      planificador domina sobre el efecto de la condición.
    \item \textbf{Los fallos no fueron de granularidad}: un timeout de agente y
      una validación cuya reparación salió de alcance.
  \end{itemize}

  \pause
  \bigskip
  \begin{center}
    \textbf{Un defecto de medición que invalida la lectura ingenua}
  \end{center}
  Los criterios de aceptación se compilan \textbf{por unidad}: descomponer
  multiplica las obligaciones. Las 12 celdas dieron cobertura 1{,}00 — pero cada
  condición satisfizo \emph{su propia} vara (5 criterios en A, 14 en B y C).
\note{%
  Este es el aporte metodologico no previsto, y conviene presentarlo como tal.
  Cualquier estudio futuro que compare estrategias de granularidad debe fijar un
  conjunto de criterios externo e identico ANTES de descomponer.%
}
\end{frame}

\begin{frame}{Lo que este diseño permite decir}
  \begin{columns}[T]
    \begin{column}{0.48\textwidth}
      \textbf{Permite}
      \begin{itemize}
        \item Caracterizar viabilidad.
        \item Describir trade-offs observados.
        \item Identificar \textbf{modos de falla}.
        \item Reportar la variabilidad del planificador.
      \end{itemize}
    \end{column}
    \begin{column}{0.48\textwidth}
      \textbf{No permite}
      \begin{itemize}
        \item Afirmar superioridad de la política.
        \item Reportar significancia.
        \item Generalizar a otros repositorios.
        \item Calibrar los pesos.
      \end{itemize}
    \end{column}
  \end{columns}

  \bigskip
  \pause
  \begin{center}
    Con $n = 2$ por celda, una muestra chica rinde en lo cualitativo.\\
    Los modos de falla son el resultado, no un subproducto.
  \end{center}
\note{%
  Esta diapositiva anticipa las preguntas incomodas. Decirlo antes de que lo
  pregunten cuesta menos que defenderlo despues. Si insisten con si sirve o no,
  la respuesta honesta: sirve para caracterizar viabilidad y modos de falla;
  para afirmar superioridad haria falta otro diseno.%
}
\end{frame}

% -----------------------------------------------------------------------------
\section{Cierre}

\begin{frame}{Conclusiones}
  \begin{itemize}
    \item Se construyó un orquestador que lleva un objetivo hasta un resultado
      \textbf{integrado, verificado y entregado}, con evidencia auditable de
      cada paso.
    \item La descomposición adaptativa está en la ruta productiva y su decisión
      queda \textbf{persistida y explicable} run por run.
    \item \textbf{Dos hallazgos negativos} son lo más transferible. Una política
      determinista puede \emph{detectar} exceso de complejidad pero no puede
      \emph{inventar el corte semántico}; y los criterios de aceptación son
      endógenos a la configuración comparada, así que toda métrica de éxito
      construida sobre ellos evalúa a cada una contra su propia vara.
    \item \textbf{No se afirma superioridad de la política adaptativa.} El
      estudio falsó su hipótesis y el documento lo reporta como tal.
    \item Los defectos que importaron no se encontraron leyendo código, sino
      \textbf{ejecutando runs reales} contra un repositorio de verdad.
  \end{itemize}
\note{%
  Cerrar por el resultado negativo, no por la plataforma. La plataforma es el
  medio; la frontera entre juicio semantico y decision determinista es lo que
  queda. Segundo mensaje: los defectos que importaron aparecieron ejecutando,
  no leyendo codigo.%
}
\end{frame}

\begin{frame}{Trabajo futuro}
  \begin{itemize}
    \item Calibrar los pesos de $C_{\mathit{task}}$ con una muestra que lo
      permita; el diseño actual, deliberadamente, no lo sostiene.
    \item Enmiendas versionadas de alcance: convertir una violación en una
      propuesta de revisión del grafo en lugar de una falla terminal.
    \item Ampliar el estudio a repositorios de dominios y tamaños distintos.
    \item Reducir la dependencia de la variabilidad del proveedor en
      planificación.
  \end{itemize}

  \vfill
  \begin{center}
    \large ¿Preguntas?
  \end{center}
\end{frame}

\end{document}
